\begin{abstract}
Reinforcement-learning (RL) detectors with asymmetric rewards have been proposed for safety-prioritized detection, where a missed event costs more than a false alarm. We ask whether such a detector, E-DDQN---a dueling Double Deep Q-Network whose reward penalizes a missed cybersickness event five times more than a false alarm---improves detection in virtual reality once evaluation is free of data leakage. On the public UFFCSData gameplay dataset (\nSessionsEval{} sessions, \nSamplesEval{} one-second samples), oversampling before a random row-level split copied \LeakDupPct\,\% of the positive test rows into the training data and gave ROC-AUC \LeakOldRocAuc; a session-grouped split of the same features gave \LeakGrpRocAuc. We prove that on logged data the alarm cannot affect future observations, so the task is a contextual bandit and the 5:1 reward induces exactly the 5:1 cost-sensitive Bayes decision rule. Under a nested, session-grouped protocol fixed before testing, E-DDQN had a mean 5:1 expected cost of \EDDQNCostCostMean{} and ROC-AUC \EDDQNCostRocAucMean. It did not outperform a standard DQN (\DQNCostCostMean, \DQNCostRocAucMean) and had a higher cost than a cost-weighted random forest (\RFCostCostMean, \RFCostRocAucMean; seed ensembles \PHEnsEDDQNCost{} vs.\ \PHEnsRFCost, Holm-adjusted $p=$\,\PHHolmRF). A post-hoc logistic regression on elapsed time alone reached ROC-AUC \PHTimeLRRocAucMean, and no model ranked seconds significantly better. At the cost-optimal thresholds every model raised alarms on \DOneFPRMin--\DOneFPRMax\,\% of non-event seconds, and bounded telemetry perturbations changed ROC-AUC by at most \AttAucMaxAbsChange. Asymmetric rewards encode a cost-sensitive operating point, not a new learning objective; RL detectors should be compared with cost-weighted and exposure-time baselines under participant-independent evaluation.
\end{abstract}

\begin{IEEEkeywords}
Cybersickness, virtual reality, deep reinforcement learning, double deep Q-network, contextual bandit, cost-sensitive learning, data leakage, safety-critical detection.
\end{IEEEkeywords}
